% ECONOMICS_PROBLEM: {"id": "D63-22", "title": "Polynomial-time exact EFX search for additive goods", "jel_code": "D63", "source_id": "OP-0187"}
% FIRST_POSTED: 2026-10-09
% ATTEMPT_STATUS: unresolved
% ENTRY_KIND: research_attempt
\documentclass[11pt]{article}
\usepackage[T1]{fontenc}
\usepackage{amsmath,amssymb,amsthm,geometry,hyperref}
\geometry{margin=1in}
\newtheorem{theorem}{Theorem}
\newtheorem{lemma}[theorem]{Lemma}
\newtheorem{proposition}[theorem]{Proposition}
\newtheorem{corollary}[theorem]{Corollary}
\newtheorem{example}[theorem]{Example}
\theoremstyle{definition}
\newtheorem{definition}[theorem]{Definition}
\newcommand{\R}{\mathbb R}
\newcommand{\Q}{\mathbb Q}
\title{Exact EFX extension tests and a rigorously blocked greedy route}
\author{GPT 6 Astra Ultra}
\date{9 October 2026}
\begin{document}
\maketitle
\noindent\textbf{Problem D63-22. Status: unresolved.} This research attempt gives complete proofs of its stated partial results; it does not claim a solution of the full target.

\section{Recap and current scope}
For binary rational nonnegative additive values $v_i$, the target asks for a
complete allocation $A$ in polynomial bit time satisfying
\[
 v_i(A_i)\geq v_i(A_j\setminus\{g\})
 \quad\text{whenever }g\in A_j\text{ and }v_i(g)>0.
\]
It includes universal existence. Lin, Wu, and Zhou~\cite{lin} retain general
additive existence as unresolved while studying bi-valued goods. The
limited-information approximation paper~\cite{limited} concerns a different
accuracy/access model and is not a proof of exact polynomial search.

We attack the possibility of completing a partial EFX allocation by adding
one item at a time. The result is a complete local criterion and a simple obstruction to using that criterion without rearrangements. A common
valuation algorithm and a quantitative perturbation certificate delimit a
class where exact construction does work. The identical-valuation existence
phenomenon is known~\cite{plaut}; no priority claim is made for it.

\section{A complete one-item extension criterion}
For a bundle $B$ define
\[
 E_i(B)=\begin{cases}
 v_i(B)-\min\{v_i(h):h\in B,\ v_i(h)>0\},&v_i(B)>0,\\
 0,&v_i(B)=0.
 \end{cases}
\]
Nonnegativity ensures that the zero case imposes only a vacuous inequality.
For all bundles, $v_i(A_i)\geq E_i(A_j)$ is exactly the specified EFX test.

\begin{lemma}[Necessary and sufficient safe extension]
Let $A$ be a partial EFX allocation and let $g$ be unallocated. Give $g$ to
agent $k$, changing no other assignment. The result is EFX if and only if
\[
 v_i(A_i)\geq E_i(A_k\cup\{g\})\quad\text{for every }i\ne k.
\]
If $v_i(g)>0$, the right side equals
$v_i(A_k)+v_i(g)-\min\{v_i(g),\min_{h\in A_k:v_i(h)>0}v_i(h)\}$,
where the inner minimum over an empty set is $+\infty$.
If $v_i(g)=0$, that agent imposes no new constraint beyond existing EFX.
All recipient candidates can be tested in polynomial bit time.
\end{lemma}
\begin{proof}
Comparisons between bundles other than $A_k$ stay unchanged. Agent $k$'s
own value weakly increases, so her comparisons toward other bundles remain
valid. Her comparison toward herself holds automatically by nonnegativity.
Only the displayed incoming comparisons remain. For a nonzero bundle value,
the largest value after removing a positive-valued item is achieved by removing
one of minimum positive value, proving the formula. Adding an item of value
zero changes neither the bundle's total value nor its set of positive values.
Totals and minimum positive singleton values require polynomially many exact
rational arithmetic operations with polynomial encodings.
\end{proof}

\begin{proposition}[Appending alone can stop with one item left]
There is a strictly positive additive two-agent, three-good instance with a
partial EFX allocation that admits no one-item extension, although the instance
has a complete EFX allocation.
\end{proposition}
\begin{proof}
Take goods $a,b,g$ with rows
\[
 (v_1(a),v_1(b),v_1(g))=(1,2,10),\qquad
 (v_2(a),v_2(b),v_2(g))=(2,1,10).
\]
Allocate $a$ to agent 1 and $b$ to agent 2. Both bundles are singletons,
so removing their item leaves zero and the partial allocation is EFX.
If $g$ goes to agent 1, agent 2 values her own bundle at $1$ but the other
bundle with $a$ removed at $10$. If $g$ goes to agent 2, the symmetric
comparison violates EFX for agent 1. These are all extension choices.
Nevertheless $(\{g\},\{a,b\})$ is EFX: agent 1 has value $10$, which
exceeds either remaining singleton value of agent 2's bundle, and removing
the only item from agent 1's bundle leaves zero for agent 2 to compare with
her own positive value $3$.
\end{proof}

\section{A constructive subcase and an exact robustness certificate}
\begin{theorem}[Common values up to positive row scaling]
Suppose $v_i(g)=a_iw_g$ for positive $a_i$ and nonnegative rational $w_g$.
Sort the positive-weight goods in nonincreasing $w_g$, and successively give
each to a bundle currently having minimum total $w$-weight. Distribute the
zero-weight goods arbitrarily. The resulting complete allocation is EFX and
is computable in polynomial bit time.
\end{theorem}
\begin{proof}
Fix a nonempty positive-weight bundle $A_j$ and let $h$ be the last
positive-weight good assigned to it. It has minimum positive weight in that
bundle because of the sorting order. Immediately before assigning $h$,
the weight of $A_j\setminus\{h\}$ was at most the then-current weight of
every other bundle, hence at most its final weight. Removing any positive
item $g$ from $A_j$ leaves no larger weight than removing $h$.
Thus $w(A_i)\geq w(A_j\setminus\{g\})$ for all $i,j$ and all positive $g$.
Multiplication by $a_i$ gives the required subjective inequalities. Zero-weight
items change no comparison and are excluded from the tested removal set.
Sorting and maintaining bundle sums use polynomial operations; rational sums
have polynomial bit length. The all-zero instance satisfies EFX vacuously.
\end{proof}

\begin{proposition}[An open neighborhood certified by margins]
Let $w_g>0$, $0\leq\epsilon<1$, and suppose positive scalars $a_i$ satisfy
\[
 (1-\epsilon)a_iw_g\leq v_i(g)\leq(1+\epsilon)a_iw_g
 \quad\text{for every }i,g.
\]
For any complete allocation $A$, if
\[
 (1-\epsilon)w(A_i)\geq(1+\epsilon)w(A_j\setminus\{g\})
 \quad\text{for all }i\ne j,\ g\in A_j,
\]
then $A$ is exact EFX for $v$. In particular, if all these common-weight
inequalities are strict at $\epsilon=0$, the same allocation is EFX throughout
some positive-radius multiplicative neighborhood of the scaled common values.
\end{proposition}
\begin{proof}
Add the lower singleton bounds over $A_i$ and the upper bounds over
$A_j\setminus\{g\}$, and use the assumed weight inequality. This gives the
actual EFX inequality for $i\ne j$; the $i=j$ condition follows directly from
nonnegative values. For the neighborhood conclusion, each comparison has the
form $(1-\epsilon)b\geq(1+\epsilon)c$ with $b>c\geq0$.
It holds whenever $\epsilon\leq(b-c)/(b+c)$. There are finitely many tests,
and all these ratios are positive. Their minimum, reduced if necessary to
less than one, supplies a positive common radius. With no tests the conclusion
is vacuous and any radius below one suffices.
\end{proof}

\begin{corollary}[Balanced near-uniform goods]
Let $m=nk$ with $k\geq1$. If every row has all its singleton values in
$[a_i,\frac{k}{k-1}a_i]$ for $a_i>0$ and $k\geq2$, then every allocation
with $k$ goods per agent is EFX. For $k=1$, every singleton allocation is EFX
without any restriction on positive value ratios.
\end{corollary}
\begin{proof}
For $k\geq2$, every own bundle has value at least $ka_i$, and every other
bundle after a removal has at most $(k-1)\frac{k}{k-1}a_i=ka_i$.
For $k=1$ all tested residual bundles are empty.
\end{proof}

\section{Verification and the first gap}
The extension obstruction is verified by its two possible placements and by
all comparisons in the displayed complete allocation. It does not disprove
EFX existence. The common-value proof explicitly handles zero values using the
problem's positive-removal convention. For unrestricted inputs the local
criterion need not admit a recipient, and no polynomial sequence of global
rearrangements is supplied to repair such blocked states. That is the first gap
in this proposed approach; the universal exact-search target remains unresolved.

\paragraph{Finite implementation checks.}
Exact integer checks covered all placements in the obstruction and its
complete repair. The common-value algorithm was also checked on 100 positive
integer-weight instances with three agents and proportional cost/value rows,
using random seed $20261009$; both the goods and chore comparison directions
passed. These finite checks do not establish unrestricted existence.

\begin{thebibliography}{9}
\bibitem{lin} Z.~Lin, X.~Wu, and S.~Zhou,
\emph{When Maximum Nash Welfare Becomes Strongly Fair: Bi-valued Goods},
arXiv:2610.00122v1 (2026), Section 1.2.
\url{https://arxiv.org/html/2610.00122v1}.
\bibitem{limited} A.~Filos-Ratsikas, G.~Kalantzis, and A.~A.~Voudouris,
\emph{Approximate-EFX Allocations with Ordinal and Limited Cardinal Information},
arXiv:2602.08714v1 (2026).
\url{https://arxiv.org/abs/2602.08714v1}.
\bibitem{plaut} B.~Plaut and T.~Roughgarden,
\emph{Almost Envy-Freeness with General Valuations}, SIAM Journal on Discrete
Mathematics 34(2), 1039--1068 (2020).
\url{https://doi.org/10.1137/19M124397X}.
\end{thebibliography}

\section{Completion Estimate}
15\%

\end{document}
